\documentclass{article}

\usepackage{amsmath,amssymb}
\usepackage{xurl}
\usepackage[hidelinks]{hyperref}
\setlength{\emergencystretch}{2em}

% Shared commands for notes in docs/paper-gaps/.

\DeclareUnicodeCharacter{2080}{\ifmmode{}_{0}\else\textsubscript{0}\fi}
\DeclareUnicodeCharacter{2081}{\ifmmode{}_{1}\else\textsubscript{1}\fi}
\DeclareUnicodeCharacter{2082}{\ifmmode{}_{2}\else\textsubscript{2}\fi}
\DeclareUnicodeCharacter{2099}{\ifmmode{}_{n}\else\textsubscript{n}\fi}

\newcommand{\Lean}{\textsc{Lean}}
\ExplSyntaxOn
\NewDocumentCommand{\leanid}{m}{
  \ifmmode
    \textsf{\detokenize{#1}}
  \else
    \group_begin:
    \urlstyle{sf}
    \tl_set:Nn \l_tmpa_tl {#1}
    \tl_replace_all:Nnn \l_tmpa_tl {\_} {_}
    \exp_args:NV \path \l_tmpa_tl
    \group_end:
  \fi
}
\ExplSyntaxOff
\providecommand{\tr}{\operatorname{tr}}
\newcommand{\tnleanIssueBase}{https://github.com/LionSR/TNLean/issues/}
\newcommand{\tnleanPrBase}{https://github.com/LionSR/TNLean/pull/}
\newcommand{\tnleanDocBase}{https://sirui-lu.com/TNLean/docs/}
\newcommand{\ghissue}[1]{\href{\tnleanIssueBase#1}{GitHub issue \##1}}
\newcommand{\ghpr}[1]{\href{\tnleanPrBase#1}{PR \##1}}
\newcommand{\doclink}[1]{\href{\tnleanDocBase#1.html}{\path{#1}}}

% Dirac notation and the identity operator, as in blueprint/src/macros/common.tex.
% \providecommand keeps note-local definitions authoritative.
\usepackage{dsfont}
\providecommand{\ket}[1]{|#1\rangle}
\providecommand{\bra}[1]{\langle #1|}
\providecommand{\braket}[2]{\langle #1 | #2 \rangle}
\providecommand{\ketbra}[2]{|#1\rangle\!\langle #2|}
\providecommand{\Id}{\mathds{1}}
\providecommand{\one}{\mathds{1}}
\providecommand{\C}{\mathbb{C}}
\providecommand{\MN}[1]{M_{#1}(\C)}
\providecommand{\MD}{M_D(\C)}

% External referencing for paper sources.  The build helper writes sanitized
% auxiliary files under build/paper-aux/ and excludes duplicated source labels.
\usepackage{xr}

% Import a generated paper auxiliary file with a caller-supplied label prefix.
% The candidate paths support builds from docs/paper-gaps/, the repository root,
% and build/paper-aux/ itself.
\newcommand{\importpaperaux}[2]{%
  \IfFileExists{../../build/paper-aux/#2.aux}{%
    \externaldocument[#1]{../../build/paper-aux/#2}%
  }{%
    \IfFileExists{build/paper-aux/#2.aux}{%
      \externaldocument[#1]{build/paper-aux/#2}%
    }{%
      \IfFileExists{#2.aux}{%
        \externaldocument[#1]{#2}%
      }{}%
    }%
  }%
}

% General external reference macros
\newcommand{\paperref}[2]{\ref{#1-#2}}
\newcommand{\papereqref}[2]{(\ref{#1-#2})}

% CPSV16 source (1606.00608 / MPDO-22-12-17-2.tex) external document and shortcuts
\importpaperaux{cpsv16-}{1606.00608/MPDO-22-12-17-2-tnlean}
\newcommand{\cpsvref}[1]{\paperref{cpsv16}{#1}}
\newcommand{\cpsveqref}[1]{\papereqref{cpsv16}{#1}}

% Verdict marker: \gapnote{<kind>}{<status>}. Kind is the policy
% classification (clarification, local-correction, scope-restriction,
% unfaithful, false-source, open-gap); status is open, wip (elimination
% underway), resolved, or historical. Machine-read by the site index;
% prints a verdict line.
\newcommand{\gapnote}[2]{\par\noindent\textbf{Verdict:} #1 (#2).\par}


\title{MPU One-Letter Mixed Residual Trace: Resolved Local Correction}
\author{}
\date{2026-08-09}

\begin{document}
\maketitle
\gapnote{local-correction}{resolved}

\section{Decomposition and source assertion}

Let $\mathcal U$ be a matrix product unitary with physical dimension $d>0$,
and let $\mathcal W$ be its physical-adjoint double layer. The source
decomposes the local tensor into identity and traceless physical components:
\begin{align}
    \mathcal W
        &=E\otimes\operatorname{Id}
          +\sum_\alpha S_\alpha\otimes\sigma_\alpha.
    \label{eq:mpu_one_letter_mixed_residual_trace_identity_traceless_decomposition}
\end{align}
Here $E$ and the residual coefficients $S_\alpha$ act on the doubled virtual
space, while each physical matrix $\sigma_\alpha$ is traceless.

For a physical matrix $X$, define its physical contraction with
$\mathcal W$ by
\begin{align}
    \operatorname{contr}_{\mathrm{phys}}(\mathcal W,X)
        &=\tr(X)E+R_{\mathcal W}(X),
    \label{eq:mpu_one_letter_mixed_residual_trace_contraction_decomposition}\\
    R_{\mathcal W}(X)
        &=\operatorname{contr}_{\mathrm{phys}}
          \left(\mathcal W,
              X-\frac{\tr(X)}{d}\operatorname{Id}\right).
    \label{eq:mpu_one_letter_mixed_residual_trace_residual_slice}
\end{align}
Thus $R_{\mathcal W}(X)$ is the residual slice associated with the traceless
part of $X$. Choosing traceless basis matrices as arguments produces the
coefficients $S_\alpha$.

The proof of Proposition~III.3(ii) of
Ref.~\cite{Cirac2017MPU} states that every positive-length closed product
containing at least one residual coefficient has zero trace
\cite[lines~405--409]{Cirac2017MPU}. The one-letter case gives
\begin{align}
    \tr(S_\alpha) &= 0.
    \label{eq:mpu_one_letter_mixed_residual_trace_single_coefficient}
\end{align}
More generally, every mixed word in $E$ and the $S_\alpha$ that contains at
least one residual letter has vanishing trace.

\section{Restricted auxiliary theorems and shifted-nilpotence recovery}

The definition of MPU unitarity directly provides closed-chain identities only
for lengths $N>1$. The intermediate constant-residual and mixed-residual
theorems therefore first establish
\begin{align}
    \tr\!\left(R_{\mathcal W}(X_0)\cdots R_{\mathcal W}(X_{N-1})\right)
        &=0,
    \qquad N>1.
    \label{eq:mpu_one_letter_mixed_residual_trace_long_residual_word}
\end{align}
The analogous result holds for mixed products of identity and residual
factors.\footnote{Formal declarations:
  \leanid{MPOTensor.IsMPU.trace_prod_residualSlice_doubleLayerTensor_eq_zero}
  and
  \leanid{MPOTensor.IsMPU.trace_prod_mixedResidualFactor_doubleLayerTensor_eq_zero}.}

The longer-word identity recovers the apparent missing one-letter case. Fix a
physical matrix $X$ and define
\begin{align}
    A &= R_{\mathcal W}(X).
    \label{eq:mpu_one_letter_mixed_residual_trace_fixed_residual_matrix}
\end{align}
Using $X$ at every site in
Eq.~\ref{eq:mpu_one_letter_mixed_residual_trace_long_residual_word} gives
\begin{align}
    \tr(A^N) &= 0,
    \qquad N>1.
    \label{eq:mpu_one_letter_mixed_residual_trace_shifted_power_traces}
\end{align}
The shifted zero-trace-power theorem implies that $A$ is nilpotent. Every
nilpotent complex matrix has zero trace, so
\begin{align}
    \tr(A) &= 0.
    \label{eq:mpu_one_letter_mixed_residual_trace_single_residual}
\end{align}
The formal one-letter theorem implements this argument.\footnote{Formal
declarations:
  \leanid{MPOTensor.IsMPU.trace_residualSlice_doubleLayerTensor_eq_zero} and
  \leanid{Matrix.isNilpotent_of_forall_trace_pow_eq_zero_of_one_lt}.}

The full positive-length theorem separates the cases $N>1$ and $N=1$. In the
one-letter branch, the unique physical site is residual, so
Eq.~\ref{eq:mpu_one_letter_mixed_residual_trace_single_residual} applies.
\footnote{Formal declaration:
  \leanid{MPOTensor.IsMPU.trace_prod_mixedResidualFactor_doubleLayerTensor_eq_zero_of_pos}.}

This proof reverses the local order of the argument in
Ref.~\cite{Cirac2017MPU}: it first proves the longer closed-word identities and
then derives the one-letter trace. It does not strengthen the MPU hypotheses.

\section{Verdict}

\textbf{Verdict: local correction (resolved).} The intermediate auxiliary
theorems retain the explicit restriction $N>1$, but shifted nilpotence derives
the missing $N=1$ case. The final positive-length mixed-residual theorem has
the same hypotheses and conclusion as the source assertion in
Ref.~\cite{Cirac2017MPU}. The former proof obligation was recorded in
\href{https://github.com/LionSR/TNLean/issues/5807}{TNLean issue~\#5807},
which closed after the derivation was formalized.

\bibliographystyle{alpha}
\bibliography{references}

\end{document}
